\documentclass[10pt,a4paper]{amsart}
\usepackage{amsmath,amssymb,amsthm,graphicx}
% For commutative diagrams we use Xy-pic
\usepackage[all]{xy}

% We use 2cell for 2-commutative diagrams.
\xyoption{2cell}
% Theorem environments.
%
\theoremstyle{plain}
\newtheorem{theorem}[subsection]{Theorem}
\newtheorem{proposition}[subsection]{Proposition}
\newtheorem{lemma}[subsection]{Lemma}

\theoremstyle{definition}
\newtheorem{definition}[subsection]{Definition}
\newtheorem{example}[subsection]{Example}
\newtheorem{exercise}[subsection]{Exercise}
\newtheorem{situation}[subsection]{Situation}

\theoremstyle{remark}
\newtheorem{remark}[subsection]{Remark}
\newtheorem{remarks}[subsection]{Remarks}

\numberwithin{equation}{subsection}

% Macros
%
\def\lim{\mathop{\mathrm{lim}}\nolimits}
\def\colim{\mathop{\mathrm{colim}}\nolimits}
\def\Spec{\mathop{\mathrm{Spec}}}
\def\Hom{\mathop{\mathrm{Hom}}\nolimits}
\def\Ext{\mathop{\mathrm{Ext}}\nolimits}
\def\SheafHom{\mathop{\mathcal{H}\!\mathit{om}}\nolimits}
\def\SheafExt{\mathop{\mathcal{E}\!\mathit{xt}}\nolimits}
\def\Sch{\mathit{Sch}}
\def\Mor{\mathop{\mathrm{Mor}}\nolimits}
\def\Ob{\mathop{\mathrm{Ob}}\nolimits}
\def\Sh{\mathop{\mathit{Sh}}\nolimits}
\def\NL{\mathop{N\!L}\nolimits}
\def\CH{\mathop{\mathrm{CH}}\nolimits}
\def\proetale{{pro\text{-}\acute{e}tale}}
\def\etale{{\acute{e}tale}}
\def\QCoh{\mathit{QCoh}}
\def\Ker{\mathop{\mathrm{Ker}}}
\def\Im{\mathop{\mathrm{Im}}}
\def\Coker{\mathop{\mathrm{Coker}}}
\def\Coim{\mathop{\mathrm{Coim}}}

% Boxtimes
%
\DeclareMathSymbol{\boxtimes}{\mathbin}{AMSa}{"02}

%
% Macros for moduli stacks/spaces
%
\def\QCohstack{\mathcal{QC}\!\mathit{oh}}
\def\Cohstack{\mathcal{C}\!\mathit{oh}}
\def\Spacesstack{\mathcal{S}\!\mathit{paces}}
\def\Quotfunctor{\mathrm{Quot}}
\def\Hilbfunctor{\mathrm{Hilb}}
\def\Curvesstack{\mathcal{C}\!\mathit{urves}}
\def\Polarizedstack{\mathcal{P}\!\mathit{olarized}}
\def\Complexesstack{\mathcal{C}\!\mathit{omplexes}}
% \Pic is the operator that assigns to X its picard group, usage \Pic(X)
% \Picardstack_{X/B} denotes the Picard stack of X over B
% \Picardfunctor_{X/B} denotes the Picard functor of X over B
\def\Pic{\mathop{\mathrm{Pic}}\nolimits}
\def\Picardstack{\mathcal{P}\!\mathit{ic}}
\def\Picardfunctor{\mathrm{Pic}}
\def\Deformationcategory{\mathcal{D}\!\mathit{ef}}
\usepackage{cleveref}
\crefname{originalreference}{original reference}{original references}
\Crefname{originalreference}{Original reference}{Original references}
\makeatletter
\newlabel{lemma-composition}{{0212}{source}{}{sourcepointer-1}{}}
\newlabel{lemma-composition@cref}{{[originalreference][0][]0212}{[1][0][]source}{}{}{}}
\makeatother
\crefname{originalreference}{original reference}{original references}
\Crefname{originalreference}{Original reference}{Original references}
\makeatletter
\newlabel{lemma-finite}{{08YX}{source}{}{sourcepointer-2}{}}
\newlabel{lemma-finite@cref}{{[originalreference][0][]08YX}{[1][0][]source}{}{}{}}
\makeatother
\crefname{originalreference}{original reference}{original references}
\Crefname{originalreference}{Original reference}{Original references}
\makeatletter
\newlabel{lemma-higher-direct-images}{{03Q8}{source}{}{sourcepointer-3}{}}
\newlabel{lemma-higher-direct-images@cref}{{[originalreference][0][]03Q8}{[1][0][]source}{}{}{}}
\makeatother
\crefname{originalreference}{original reference}{original references}
\Crefname{originalreference}{Original reference}{Original references}
\makeatletter
\newlabel{lemma-sandwich}{{0A4C}{source}{}{sourcepointer-4}{}}
\newlabel{lemma-sandwich@cref}{{[originalreference][0][]0A4C}{[1][0][]source}{}{}{}}
\makeatother
\crefname{originallemma}{original lemma}{original lemmas}
\Crefname{originallemma}{Original lemma}{Original lemmas}
\makeatletter
\newlabel{limits-lemma-chow-finite-type}{{0202}{source}{}{sourcepointer-5}{}}
\newlabel{limits-lemma-chow-finite-type@cref}{{[originallemma][0][]0202}{[1][0][]source}{}{}{}}
\makeatother
\crefname{originallemma}{original lemma}{original lemmas}
\Crefname{originallemma}{Original lemma}{Original lemmas}
\makeatletter
\newlabel{morphisms-lemma-composition-proper}{{01W3}{source}{}{sourcepointer-6}{}}
\newlabel{morphisms-lemma-composition-proper@cref}{{[originallemma][0][]01W3}{[1][0][]source}{}{}{}}
\makeatother
\crefname{originallemma}{original lemma}{original lemmas}
\Crefname{originallemma}{Original lemma}{Original lemmas}
\makeatletter
\newlabel{morphisms-lemma-image-proper-scheme-closed}{{01W6}{source}{}{sourcepointer-7}{}}
\newlabel{morphisms-lemma-image-proper-scheme-closed@cref}{{[originallemma][0][]01W6}{[1][0][]source}{}{}{}}
\makeatother
\title{Base change for all proper morphisms reduces to the projective line over every base}
\author{The Stacks Project Authors}
\date{stacks-project:etale-cohomology.tex:local-snapshot-20261011; GNU FDL 1.2 or later}
\begin{document}
\maketitle
\noindent\textit{Editorial note.} All delivered mathematical text is the author's own, copied from the pinned Stacks Project chapter file etale-cohomology.tex; no statement or proof step was rewritten, reordered or removed.\par\medskip
\noindent\textit{Editorial note.} The counted proof is the author's complete argument: the proof is written out in full: the localisation to an affine base, the Chow's lemma diagram with the closed immersion, and the two rounds of the sandwich/finite/composition reductions ending with the polynomial factorisation of the projection.\par\medskip
\noindent\textit{Editorial note.} The label lemma-reduce-to-P1 names the counted statement, and the proof environment immediately following it in the source is delivered verbatim as its proof.\par\medskip
\noindent\textit{Editorial note.} The counted statement is carried by the author's own lemma environment in the chapter source.\par\medskip
\noindent\textit{Editorial note.} The statement and proof are copied from the pinned chapter file; the ten cited inputs are kept as external pointers with their canonical Stacks tags.\par\medskip
\noindent\textit{Allowed standard prerequisites (editorial).}\begin{itemize}
\item Formation of higher direct images commutes with localisation on the base. Same chapter, the lemma stating that the higher direct images may be computed after restricting to an open of the base. Supplies the reduction to an affine base.
\item Chow's lemma and the geography of the resulting diagram. Limits, the form of Chow's lemma giving X' -> X proper surjective with X' -> P\textasciicircum{}n\_Y an immersion, and Morphisms, the lemmas that a composite of proper morphisms is proper and that a proper morphism with an immersion into a proper scheme is a closed immersion. Supplies the commutative diagram with the closed immersion X' -> P\textasciicircum{}n\_X.
\item The sandwich, finite and composition principles. Same chapter, the lemma deducing base change for a composite from the two factors, the lemma for finite morphisms and the lemma for composites. Supplies the reductions used twice in the proof.
\end{itemize}\medskip
\ifdefined\hypertarget\hypertarget{sourcepointer-1}{}\fi
\noindent\textit{Source reference.} Original reference 0212: sources/stacks-etale-cohomology/pinned-original-source.tex, source lines 18477--18477. Prerequisite pointer invoked at proof line(s) 18583; it names a standard fact of the same chapter that the counted argument uses, and every step of the counted proof built on it is retained verbatim in the delivered body.\par\medskip
\ifdefined\hypertarget\hypertarget{sourcepointer-2}{}\fi
\noindent\textit{Source reference.} Original reference 08YX: sources/stacks-etale-cohomology/pinned-original-source.tex, source lines 18519--18519. Prerequisite pointer invoked at proof line(s) 18581; it names a standard fact of the same chapter that the counted argument uses, and every step of the counted proof built on it is retained verbatim in the delivered body.\par\medskip
\ifdefined\hypertarget\hypertarget{sourcepointer-3}{}\fi
\noindent\textit{Source reference.} Original reference 03Q8: sources/stacks-etale-cohomology/pinned-original-source.tex, source lines 7510--7510. Prerequisite pointer invoked at proof line(s) 18554; it names a standard fact of the same chapter that the counted argument uses, and every step of the counted proof built on it is retained verbatim in the delivered body.\par\medskip
\ifdefined\hypertarget\hypertarget{sourcepointer-4}{}\fi
\noindent\textit{Source reference.} Original reference 0A4C: sources/stacks-etale-cohomology/pinned-original-source.tex, source lines 18431--18431. Prerequisite pointer invoked at proof line(s) 18577; it names a standard fact of the same chapter that the counted argument uses, and every step of the counted proof built on it is retained verbatim in the delivered body.\par\medskip
\ifdefined\hypertarget\hypertarget{sourcepointer-5}{}\fi
\noindent\textit{Source reference.} Original lemma 0202: sources/stacks-chow/tags.txt, source lines 2619--2619. Prerequisite pointer invoked at proof line(s) 18568; it names a standard fact of a sibling chapter that the counted argument uses, and every step of the counted proof built on it is retained verbatim in the delivered body.\par\medskip
\ifdefined\hypertarget\hypertarget{sourcepointer-6}{}\fi
\noindent\textit{Source reference.} Original lemma 01W3: sources/stacks-chow/tags.txt, source lines 2468--2468. Prerequisite pointer invoked at proof line(s) 18570; it names a standard fact of a sibling chapter that the counted argument uses, and every step of the counted proof built on it is retained verbatim in the delivered body.\par\medskip
\ifdefined\hypertarget\hypertarget{sourcepointer-7}{}\fi
\noindent\textit{Source reference.} Original lemma 01W6: sources/stacks-chow/tags.txt, source lines 2471--2471. Prerequisite pointer invoked at proof line(s) 18572; it names a standard fact of a sibling chapter that the counted argument uses, and every step of the counted proof built on it is retained verbatim in the delivered body.\par\medskip
\begin{lemma}
\label{lemma-reduce-to-P1}
To prove that cohomology commutes with base change for
every proper morphism of schemes it suffices to prove it
holds for the morphism $\mathbf{P}^1_S \to S$ for every scheme $S$.
\end{lemma}

\begin{proof}
Let $f : X \to Y$ be a proper morphism of schemes.
Let $Y = \bigcup Y_i$ be an affine open covering
and set $X_i = f^{-1}(Y_i)$. If we can prove
cohomology commutes with base change for $X_i \to Y_i$,
then cohomology commutes with base change for $f$.
Namely, the formation of the higher direct images
commutes with Zariski (and even \'etale) localization
on the base, see
Lemma \ref{lemma-higher-direct-images}.
Thus we may assume $Y$ is affine.

\medskip\noindent
Let $Y$ be an affine scheme and let $X \to Y$ be a proper morphism.
By Chow's lemma there exists a commutative diagram
$$
\xymatrix{
X \ar[rd] & X' \ar[d] \ar[l]^\pi \ar[r] & \mathbf{P}^n_Y \ar[dl] \\
& Y &
}
$$
where $X' \to \mathbf{P}^n_Y$ is an immersion, and
$\pi : X' \to X$ is proper and surjective, see
Limits, Lemma \ref{limits-lemma-chow-finite-type}.
Since $X \to Y$ is proper, we find that $X' \to Y$ is proper
(Morphisms, Lemma \ref{morphisms-lemma-composition-proper}).
Hence $X' \to \mathbf{P}^n_Y$ is a closed immersion
(Morphisms, Lemma \ref{morphisms-lemma-image-proper-scheme-closed}).
It follows that $X' \to X \times_Y \mathbf{P}^n_Y = \mathbf{P}^n_X$
is a closed immersion (as an immersion with closed image).

\medskip\noindent
By Lemma \ref{lemma-sandwich}
it suffices to prove cohomology commutes with base change for
$\pi$ and $X' \to Y$. These morphisms both factor as a closed
immersion followed by a projection $\mathbf{P}^n_S \to S$ (for some $S$).
By Lemma \ref{lemma-finite} the result holds for closed
immersions (as closed immersions are finite).
By Lemma \ref{lemma-composition} it suffices to prove the
result for projections $\mathbf{P}^n_S \to S$.

\medskip\noindent
For every $n \geq 1$ there is a finite surjective morphism
$$
\mathbf{P}^1_S \times_S \ldots \times_S \mathbf{P}^1_S
\longrightarrow
\mathbf{P}^n_S
$$
given on coordinates by
$$
((x_1 : y_1), (x_2 : y_2), \ldots, (x_n : y_n))
\longmapsto
(F_0 : \ldots : F_n)
$$
where $F_0, \ldots, F_n$ in $x_1, \ldots, y_n$
are the polynomials with integer coefficients such that
$$
\prod (x_i t + y_i) = F_0 t^n + F_1 t^{n - 1} + \ldots + F_n
$$
Applying
Lemmas \ref{lemma-sandwich}, \ref{lemma-finite}, and \ref{lemma-composition}
one more time we conclude that the lemma is true.
\end{proof}

\end{document}
